%-------------------------------------------------------------------------------
%	SECTION TITLE
%-------------------------------------------------------------------------------
\cvsection{Research Experience}

\begin{cvparagraph}
    \textbf{{\fontsize{10}{12}\selectfont Research Interests: }} \\
    Machine learning–based surrogate modeling, multi-fidelity analysis, uncertainty quantification, and design automation. I integrate physics-informed ML surrogates, multidisciplinary solvers, and automated workflows into new product development environment to accelerate assessment, multidisciplinary design optimization (MDO), and design space exploration to enable the development of sustainable unconventional vehicles with reduced program risk.
\end{cvparagraph}
%-------------------------------------------------------------------------------
%	CONTENT
%-------------------------------------------------------------------------------
\begin{cventries}

%---------------------------------------------------------
  \cventry
    {Postdoctoral Researcher} % Job title
    {University of Washington, NSF AI Institute in Dynamic Systems} % Organization
    {Seattle, WA} % Location
    {Oct. 2026 - now} % Date(s)
    {\begin{cvitems}
    \item {Developing state-of-the-art benchmarks for coupled aerodynamic, structural, and control optimization..}
    \end{cvitems}
    }

%---------------------------------------------------------
  \cventry
    {Postdoctoral Researcher} % Job title
    {Stanford University, Aerospace Design Lab} % Organization
    {Palo Alto, CA} % Location
    {Jul. 2026 - Sep. 2026} % Date(s)
    {\begin{cvitems}
    \item {Metroplex Overflight Noise Analysis (MONA) departure noise analyst.}
       \end{cvitems}
    }



%---------------------------------------------------------
  \cventry
    {Doctoral Researcher} % Job title
    {Stanford University, Aerospace Design Lab} % Organization
    {Palo Alto, CA} % Location
    {Jan. 2022 - Jun. 2026} % Date(s)
    {
      \begin{cvitems} % Description(s) of tasks/responsibilities
        \item {Advisor and Lab PI: Juan J. Alonso}
        \item {Research focuses on physics‑informed machine learning, applied surrogate modeling, AI driven design space exploration, automated design tools, and multi‑fidelity analysis methods for sustainable unconventional vehicle developments.}
        \item {Developed VortexNet, a first-of-its-kind physics-informed GNN that augments low-fidelity aerodynamic tools with high-fidelity CFD field data for conceptual aircraft design.}
        \item {Established empirical data-scaling law for VortexNet and released a multi-fidelity aerodynamic dataset.}
        \item {Demonstrated VortexNet’s integration into the design pipeline through a stability-constrained shape optimization problem.}
        \item {Contributor of the Supersonic Configurations at Low Speeds (SCALOS) project with specialization in surrogate modeling, SST design space exploration, and planform optimization.}
        \item {6 publications in the domain of applied surrogate modeling using machine learning (Physics AI), 1 publication in the domain of agentic engineering AI, and 1 publication in the domain of vehicle design-analysis tool.}
      \end{cvitems}
    }

% --------------
\cventry
    {Research Fellow} % Job title
    {GE Aerospace Research, Probabilistic Design and Material Informatics} % Organization
    {Niskayuna, NY} % Location
    {June 2025 - Sep. 2025} % Date(s)
    {
      \begin{cvitems} % Description(s) of tasks/responsibilities
        \item {Proposed and composed a Large Eddy Simulation (LES) flow super-resolution program, that enhances low-order LES flow field to high-order equivalent with sufficient credibility in integrated quantities of interests statistics.}
        \item {Surveyed the state-of-the-art methods in physics-informed machine learning and integrated selected methods into GE property software environment.}
      \end{cvitems}
    }
%---------------------------------------------------------
  \cventry
    {Research Intern} % Job title
    {SRI International, Intelligent Systems Lab} % Organization
    {Palo Alto, CA} % Location
    {Jun. 2023 - Sep. 2023} % Date(s)
    {
      \begin{cvitems} % Description(s) of tasks/responsibilities
        \item {Composed one Finite Element Analysis (FEA) program for structural simulation using MITC4 Shell element.}
        \item {Integrated the FEA code with topology optimization heuristic to enable reinforcement fiber layout optimization for engineering design optimization applications.}
      \end{cvitems}
    }

%---------------------------------------------------------
  \cventry
    {Visiting Scholar} % Job title
    {Zhejiang University, Lab of Flexible Sensors and Intelligent Interaction} % Organization
    {Hangzhou, China} % Location
    {Aug. 2020 - October. 2021} % Date(s)
    {
      \begin{cvitems} % Description(s) of tasks/responsibilities
        \item {Performed conformal mapping simulation for hyper-elastic soft sensor deployment on complex robot arm surface.}
        \item {Proposed, designed, and simulated one soft hair-like flow sensor that is applicable for small drone air speed sensing.}
      \end{cvitems}
    }

%---------------------------------------------------------
  \cventry
    {Research Assistant} % Job title
    {Cornell University, Thermal Science Lab} % Organization
    {Ithaca, NY} % Location
    {Jan. 2015 - Jan. 2017} % Date(s)
    {
      \begin{cvitems} % Description(s) of tasks/responsibilities
        \item {Advisor and Lab PI: C. Thomas Avedisian Ph.D.}
        \item {Research focuses on fuel droplet combustion assessment, including burning rate and soot production, for surrogate gasoline fuel, jet fuel and renewable jet fuel. Developed various automatic image analysis software to facilitate data acquisition.}
      \end{cvitems}
    }

\end{cventries}
